\documentclass[12pt, a4paper]{article}

% ─── Encoding & Language ───────────────────────────────────────────────────────
\usepackage[utf8]{inputenc}
\usepackage[T1]{fontenc}
\usepackage[italian]{babel}

% ─── Typography ────────────────────────────────────────────────────────────────
\usepackage{lmodern}
\usepackage{microtype}
\usepackage{setspace}
\onehalfspacing

% ─── Page Layout ───────────────────────────────────────────────────────────────
\usepackage[
  top=2.5cm, bottom=2.5cm,
  left=3cm,  right=2.5cm
]{geometry}

% ─── Colors ────────────────────────────────────────────────────────────────────
\usepackage[dvipsnames, table]{xcolor}
\definecolor{primary}{HTML}{2563EB}
\definecolor{codebg}{HTML}{F3F4F6}
\definecolor{codefg}{HTML}{1F2937}
\definecolor{gps}{HTML}{166534}
\definecolor{network}{HTML}{1E3A8A}
\definecolor{ipcolor}{HTML}{78350F}

% ─── Code Listings ─────────────────────────────────────────────────────────────
\usepackage{listings}
\lstdefinelanguage{TypeScript}{
  keywords={
    import, export, from, const, let, var, function, async, await,
    return, if, else, try, catch, throw, new, typeof, null, true, false,
    type, interface, enum, extends, implements, class, void
  },
  keywordstyle=\color{primary}\bfseries,
  ndkeywords={string, number, boolean, Record, Promise},
  ndkeywordstyle=\color{MidnightBlue},
  sensitive=true,
  comment=[l]{//},
  morecomment=[s]{/*}{*/},
  commentstyle=\color{gray}\itshape,
  stringstyle=\color{ForestGreen},
  morestring=[b]",
  morestring=[b]',
  morestring=[b]`
}
\lstset{
  language=TypeScript,
  backgroundcolor=\color{codebg},
  basicstyle=\ttfamily\small\color{codefg},
  numbers=left,
  numberstyle=\tiny\color{gray},
  numbersep=8pt,
  stepnumber=1,
  breaklines=true,
  breakatwhitespace=true,
  frame=single,
  framerule=0.4pt,
  rulecolor=\color{gray!40},
  xleftmargin=1.5em,
  framexleftmargin=1.5em,
  showstringspaces=false,
  tabsize=2,
  captionpos=b,
  aboveskip=1em,
  belowskip=0.5em
}

% ─── Tables ────────────────────────────────────────────────────────────────────
\usepackage{booktabs}
\usepackage{tabularx}
\usepackage{array}
\newcolumntype{L}[1]{>{\raggedright\arraybackslash}p{#1}}
\newcolumntype{C}[1]{>{\centering\arraybackslash}p{#1}}

% ─── Diagrams & Graphics ───────────────────────────────────────────────────────
\usepackage{tikz}
\usetikzlibrary{shapes.geometric, arrows.meta, positioning, fit, backgrounds}
\usepackage{graphicx}

% ─── Hyperlinks ────────────────────────────────────────────────────────────────
\usepackage[
  colorlinks=true,
  linkcolor=primary,
  urlcolor=primary,
  citecolor=primary
]{hyperref}

% ─── Miscellaneous ─────────────────────────────────────────────────────────────
\usepackage{enumitem}
\usepackage{fancyvrb}
\usepackage{tcolorbox}
\tcbuselibrary{skins, breakable}

% ─── Custom Boxes ──────────────────────────────────────────────────────────────
\newtcolorbox{notebox}{
  enhanced, breakable,
  colback=blue!5, colframe=primary,
  arc=4pt, boxrule=0.6pt,
  left=8pt, right=8pt, top=6pt, bottom=6pt,
  fonttitle=\bfseries, title={Nota}
}
\newtcolorbox{warningbox}{
  enhanced, breakable,
  colback=orange!8, colframe=orange!70!black,
  arc=4pt, boxrule=0.6pt,
  left=8pt, right=8pt, top=6pt, bottom=6pt,
  fonttitle=\bfseries, title={Attenzione}
}

% ─── Header / Footer ───────────────────────────────────────────────────────────
\usepackage{fancyhdr}
\pagestyle{fancy}
\fancyhf{}
\fancyhead[L]{\small\textit{GameInApp — Documentazione Tecnica}}
\fancyhead[R]{\small\textit{Geolocalizzazione}}
\fancyfoot[C]{\thepage}
\renewcommand{\headrulewidth}{0.4pt}

% ═══════════════════════════════════════════════════════════════════════════════
\begin{document}
% ═══════════════════════════════════════════════════════════════════════════════

\begin{titlepage}
  \centering
  \vspace*{3cm}
  {\huge\bfseries GameInApp\par}
  \vspace{0.5cm}
  {\Large Documentazione Tecnica\par}
  \vspace{1.5cm}
  \rule{\linewidth}{0.5pt}\\[0.4cm]
  {\LARGE\bfseries\color{primary} Sistema di Geolocalizzazione\par}
  \rule{\linewidth}{0.5pt}
  \vspace{2cm}
  {\large Piera Di Fusco\par}
  \vspace{0.5cm}
  {\normalsize Tesi di Laurea — Ingegneria Informatica\par}
  \vfill
  {\large \today\par}
\end{titlepage}

\tableofcontents
\newpage

% ═══════════════════════════════════════════════════════════════════════════════
\section{Introduzione}
% ═══════════════════════════════════════════════════════════════════════════════

Il sistema di geolocalizzazione di \textbf{GameInApp} è una delle funzionalità
core dell'applicazione: determina la posizione geografica dell'utente per
restituire eventi sportivi e palestre nelle immediate vicinanze.

La progettazione di questa funzionalità ha richiesto particolare attenzione
alla \textbf{robustezza} e alla \textbf{degradazione controllata}: il sistema
deve funzionare correttamente in una grande varietà di scenari, tra cui
dispositivi mobili con GPS hardware, computer desktop privi di sensore di
posizione, browser con autorizzazioni negate dall'utente, e ambienti di rete
con connettività limitata.

\subsection{Obiettivi di progettazione}

\begin{itemize}[leftmargin=1.5em]
  \item \textbf{Massima copertura}: la posizione deve essere determinabile
    in qualsiasi scenario ragionevole, con fallback graduali.
  \item \textbf{Trasparenza all'utente}: il sistema comunica sempre il livello
    di precisione ottenuto tramite un indicatore visivo.
  \item \textbf{Manutenibilità}: la logica di geolocalizzazione è incapsulata
    in un custom React Hook riutilizzabile, separata dalla logica di
    presentazione.
  \item \textbf{Sicurezza}: le dipendenze da servizi terzi sono esposte solo
    lato server, mai nel bundle client.
\end{itemize}

% ═══════════════════════════════════════════════════════════════════════════════
\section{Architettura del sistema}
% ═══════════════════════════════════════════════════════════════════════════════

\subsection{Panoramica dei componenti}

Il sistema si compone di tre elementi principali, ciascuno con una
responsabilità ben definita:

\begin{table}[h!]
  \centering
  \caption{Componenti del sistema di geolocalizzazione}
  \begin{tabularx}{\linewidth}{L{4.5cm} C{2.5cm} L{6.5cm}}
    \toprule
    \textbf{File} & \textbf{Tipo} & \textbf{Responsabilità} \\
    \midrule
    \texttt{lib/hooks/useGeolocation.ts} & React Hook & Logica client-side, strategia a livelli, gestione stato \\
    \addlinespace
    \texttt{app/api/geo/route.ts} & API Route & Proxy server-side per IP geolocation \\
    \addlinespace
    \texttt{app/page.tsx} & Componente & Consumo del hook, badge indicatore source \\
    \bottomrule
  \end{tabularx}
\end{table}

\subsection{Strategia a tre livelli}

Il sistema adotta una strategia a cascata (\textit{cascade fallback}) con tre
livelli di precisione decrescente. Se un livello fallisce, si passa
automaticamente al successivo senza alcun intervento dell'utente.

\vspace{0.8em}
\begin{center}
\begin{tikzpicture}[
  node distance=0.6cm and 0cm,
  box/.style={
    rectangle, rounded corners=4pt,
    text width=7cm, align=center,
    minimum height=1cm, font=\small,
    draw, thick
  },
  arrow/.style={-{Stealth[length=6pt]}, thick, gray},
  label/.style={font=\footnotesize\itshape, text=gray}
]

\node[box, fill=green!12, draw=green!60!black]     (l1)
  {\textbf{Livello 1 — GPS ad alta precisione}\\
   \small\texttt{enableHighAccuracy: true} — timeout 8\,s};

\node[box, fill=blue!10, draw=blue!50!black,
      below=1.4cm of l1]                            (l2)
  {\textbf{Livello 2 — Geolocalizzazione di rete}\\
   \small\texttt{enableHighAccuracy: false} — timeout 8\,s};

\node[box, fill=orange!10, draw=orange!60!black,
      below=1.4cm of l2]                            (l3)
  {\textbf{Livello 3 — IP Geolocation (server-side)}\\
   \small\texttt{/api/geo} → ip-api.com / ipapi.co};

\node[box, fill=red!10, draw=red!50!black,
      below=1.4cm of l3]                            (err)
  {\textbf{Errore definitivo}\\
   \small Messaggio localizzato + bottone \textit{Riprova}};

\draw[arrow] (l1) -- node[right=4pt, label]
  {code 2 (UNAVAILABLE) o code 3 (TIMEOUT)} (l2);
\draw[arrow] (l2) -- node[right=4pt, label]
  {qualsiasi errore} (l3);
\draw[arrow] (l3) -- node[right=4pt, label]
  {entrambi i provider falliscono} (err);

\end{tikzpicture}
\end{center}

% ═══════════════════════════════════════════════════════════════════════════════
\section{Custom Hook: \texttt{useGeolocation}}
% ═══════════════════════════════════════════════════════════════════════════════

\subsection{Interfaccia pubblica}

Il hook espone la seguente interfaccia TypeScript:

\begin{lstlisting}[caption={Tipo di ritorno del hook \texttt{useGeolocation}}]
type GeoCoords = { lat: number; lng: number };
type GeoSource = "gps" | "network" | "ip" | null;

function useGeolocation(): {
  coords:  GeoCoords | null; // coordinate o null se non ancora note
  source:  GeoSource;        // sorgente della posizione
  loading: boolean;          // true durante il processo di rilevamento
  error:   string | null;    // messaggio d'errore in italiano
  retry:   () => void;       // riavvia il processo senza ricaricare la pagina
}
\end{lstlisting}

\subsection{Significato di \texttt{source}}

Il campo \texttt{source} comunica all'utente la precisione della posizione
rilevata, permettendo di mostrare un indicatore visivo appropriato:

\begin{table}[h!]
  \centering
  \caption{Valori del campo \texttt{source} e precisione associata}
  \begin{tabularx}{\linewidth}{C{2cm} L{5cm} L{5.5cm}}
    \toprule
    \textbf{Valore} & \textbf{Sorgente} & \textbf{Precisione tipica} \\
    \midrule
    \texttt{"gps"}     & GPS hardware o WiFi Positioning ad alta accuratezza & $< 10$\,m \\
    \addlinespace
    \texttt{"network"} & Geolocalizzazione browser a bassa accuratezza (WiFi/IP lato browser) & $100$\,m -- $5$\,km \\
    \addlinespace
    \texttt{"ip"}      & Risoluzione tramite indirizzo IP pubblico & Livello città ($\sim 10$\,km) \\
    \addlinespace
    \texttt{null}      & Posizione non ancora determinata (\textit{loading}) & — \\
    \bottomrule
  \end{tabularx}
\end{table}

\subsection{Gestione degli errori W3C}

La Geolocation API del browser definisce tre codici di errore standardizzati
\cite{w3c-geolocation}. Il sistema li gestisce in modo differenziato:

\begin{table}[h!]
  \centering
  \caption{Gestione dei codici errore della Geolocation API W3C}
  \begin{tabularx}{\linewidth}{C{1.2cm} L{3.8cm} L{7.5cm}}
    \toprule
    \textbf{Code} & \textbf{Costante} & \textbf{Comportamento} \\
    \midrule
    \texttt{1} & \texttt{PERMISSION\_DENIED}      & Errore immediato, nessun retry. L'utente ha esplicitamente negato l'accesso alla posizione. \\
    \addlinespace
    \texttt{2} & \texttt{POSITION\_UNAVAILABLE}   & Retry a bassa accuratezza (Livello 2) → fallback IP (Livello 3). \\
    \addlinespace
    \texttt{3} & \texttt{TIMEOUT}                 & Retry a bassa accuratezza (Livello 2) → fallback IP (Livello 3). \\
    \addlinespace
    —          & API non disponibile              & Salta direttamente al fallback IP (Livello 3). \\
    \bottomrule
  \end{tabularx}
\end{table}

\begin{notebox}
Il codice \texttt{1 (PERMISSION\_DENIED)} non viene mai ritentato
automaticamente: una volta che l'utente ha negato il permesso, ritentare
sarebbe inutile e potenzialmente fastidioso. Viene invece mostrato un bottone
\textit{Riprova} che l'utente può premere \emph{dopo} aver modificato le
impostazioni del browser.
\end{notebox}

\subsection{Implementazione}

\subsubsection{Cancellazione e prevenzione dei memory leak}

Un problema comune con i custom hook asincroni in React è l'aggiornamento di
stato su componenti già smontati, che provoca warning e potenziali memory
leak. Il hook usa il pattern \texttt{cancelled flag} per prevenirlo:

\begin{lstlisting}[caption={Pattern cancelled flag per cleanup asincrono}]
useEffect(() => {
  let cancelled = false;

  const succeed = (lat: number, lng: number, src: GeoSource) => {
    if (cancelled) return; // no setState su componente smontato
    setCoords({ lat, lng });
    setSource(src);
    setLoading(false);
  };

  // ... logica geolocalizzazione ...

  return () => {
    cancelled = true; // cleanup: le callback ignoreranno i risultati
  };
}, [attempt]);
\end{lstlisting}

\subsubsection{Meccanismo di retry}

La funzione \texttt{retry()} sfrutta il pattern idiomatico React di
incrementare un contatore come dipendenza del \texttt{useEffect}, evitando
la necessità di logica imperativa:

\begin{lstlisting}[caption={Implementazione del retry tramite contatore}]
const [attempt, setAttempt] = useState(0);

const retry = useCallback(() => {
  setCoords(null);
  setSource(null);
  setError(null);
  setLoading(true);
  setAttempt((n) => n + 1); // incremento -> re-run del useEffect
}, []);

useEffect(() => {
  // ... tutto il processo di geolocalizzazione ...
}, [attempt]); // re-eseguito ad ogni chiamata di retry()
\end{lstlisting}

% ═══════════════════════════════════════════════════════════════════════════════
\section{API Route: \texttt{GET /api/geo}}
% ═══════════════════════════════════════════════════════════════════════════════

\subsection{Motivazione del proxy server-side}

La geolocalizzazione tramite IP potrebbe sembrare una funzionalità da
implementare direttamente nel browser (chiamata client-side al provider).
Tuttavia, l'adozione di un proxy server-side offre vantaggi significativi:

\begin{enumerate}[leftmargin=1.5em]
  \item \textbf{Nessun problema di CORS}: le chiamate avvengono
    server $\rightarrow$ provider, non browser $\rightarrow$ provider.
  \item \textbf{Nessuna dipendenza nel bundle client}: il provider può essere
    cambiato senza modificare o ridistribuire il codice frontend.
  \item \textbf{Caching server-side}: la risposta viene cachata per 5 minuti
    (\texttt{Cache-Control} e \texttt{next: \{ revalidate: 300 \}}), riducendo
    le chiamate al provider.
  \item \textbf{Sicurezza}: un'eventuale API key futura non sarà mai esposta
    nel codice client.
\end{enumerate}

\subsection{Catena di provider}

L'API route implementa una catena di provider con failover automatico:

\begin{lstlisting}[caption={Catena di provider in \texttt{/api/geo}}]
// Provider 1: ip-api.com
// - Gratuito, 1000 req/min, nessuna API key
// - Fallisce con IP privati (es. 127.0.0.1 in sviluppo locale)
const ipApiRes = await fetch(
  `http://ip-api.com/json/${useIp}?fields=status,lat,lon,city,country`
);

// Provider 2: ipapi.co (fallback)
// - Gratuito (~1000 req/giorno), funziona anche da localhost
const ipapiRes = await fetch("https://ipapi.co/json/");
\end{lstlisting}

\subsection{Gestione degli indirizzi IP privati}

In ambiente di sviluppo locale, l'header \texttt{X-Forwarded-For} non è
presente (nessun reverse proxy davanti a Next.js), quindi il server riceve
l'IP di loopback \texttt{127.0.0.1} o \texttt{::1}. Questi indirizzi
appartengono a range privati RFC 1918 \cite{rfc1918} che i provider di
geolocalizzazione non sono in grado di risolvere.

La funzione \texttt{isPrivateIp()} rileva questi casi:

\begin{lstlisting}[caption={Rilevamento di indirizzi IP privati}]
function isPrivateIp(ip: string): boolean {
  return (
    ip === "127.0.0.1" || ip === "::1" || ip === "localhost" ||
    /^10\./.test(ip) ||                   // RFC 1918 Class A
    /^172\.(1[6-9]|2\d|3[01])\./.test(ip) || // RFC 1918 Class B
    /^192\.168\./.test(ip) ||             // RFC 1918 Class C
    /^::ffff:127\./.test(ip)              // IPv4-mapped loopback
  );
}
\end{lstlisting}

\subsection{Formato della risposta}

In caso di successo, l'endpoint restituisce:

\begin{lstlisting}[language={},caption={Risposta JSON di \texttt{GET /api/geo}}]
{
  "lat":      41.9028,
  "lng":      12.4964,
  "city":     "Roma",
  "country":  "Italy",
  "accuracy": "city"
}
\end{lstlisting}

Il campo \texttt{accuracy} indica esplicitamente che si tratta di una
posizione a livello di città, permettendo all'interfaccia di mostrare il
badge \textit{IP (approx.)} appropriato.

% ═══════════════════════════════════════════════════════════════════════════════
\section{Integrazione nell'interfaccia utente}
% ═══════════════════════════════════════════════════════════════════════════════

\subsection{Consumo del hook}

In \texttt{app/page.tsx}, tutta la logica di geolocalizzazione è ridotta a
una singola riga di destructuring:

\begin{lstlisting}[caption={Consumo del hook nella pagina principale}]
const {
  coords,            // { lat, lng } passato a /api/events/search e alla mappa
  source: geoSource, // "gps" | "network" | "ip" | null
  loading: geoLoading,
  error: geoError,
  retry: retryGeo,   // collegato al bottone "Riprova"
} = useGeolocation();

// Lo stato di loading e error viene composto con quelli del fetch eventi
const loading = geoLoading || fetchLoading;
const error   = geoError ?? fetchError;
\end{lstlisting}

\subsection{Indicatore visivo della sorgente}

Per massima trasparenza verso l'utente, viene mostrato un badge colorato che
indica la qualità della posizione rilevata:

\begin{table}[h!]
  \centering
  \caption{Badge visivi per la sorgente di posizione}
  \begin{tabularx}{\linewidth}{C{2cm} C{2.5cm} C{2.5cm} L{5cm}}
    \toprule
    \textbf{Source} & \textbf{Label} & \textbf{Colore} & \textbf{Messaggio implicito} \\
    \midrule
    \texttt{"gps"}     & GPS           & Verde  & Posizione precisa \\
    \texttt{"network"} & Rete          & Blu    & Posizione approssimativa \\
    \texttt{"ip"}      & IP (approx.)  & Ambra  & Solo città — possibile imprecisione \\
    \bottomrule
  \end{tabularx}
\end{table}

% ═══════════════════════════════════════════════════════════════════════════════
\section{Limitazioni e considerazioni}
% ═══════════════════════════════════════════════════════════════════════════════

\subsection{Precisione della geolocalizzazione IP}

La geolocalizzazione tramite IP è intrinsecamente limitata nella precisione.
I provider mappano gli indirizzi IP a blocchi assegnati dagli ISP, che
possono corrispondere a regioni geografiche anche molto ampie. In ambiente
urbano la precisione è solitamente a livello di città (qualche chilometro),
ma in aree rurali o con ISP regionali può essere molto peggiore.

\begin{warningbox}
L'utente che utilizza una VPN vedrà una posizione corrispondente al nodo VPN,
non alla sua posizione reale. Il badge \textit{IP (approx.)} aiuta a
comunicare questa incertezza.
\end{warningbox}

\subsection{Privacy e GDPR}

La posizione geografica è un dato personale ai sensi del
Regolamento~(UE)~2016/679 (GDPR)~\cite{gdpr}. Il sistema:

\begin{itemize}[leftmargin=1.5em]
  \item Richiede il consenso esplicito dell'utente tramite il popup nativo del
    browser (livelli 1 e 2);
  \item Non persiste le coordinate in nessun database;
  \item Usa le coordinate esclusivamente per la query di ricerca eventi
    (\texttt{/api/events/search});
  \item Nel caso del fallback IP, non trasmette l'IP al client — il proxy
    server-side effettua la chiamata al provider esterno.
\end{itemize}

\subsection{Rate limiting dei provider}

\begin{table}[h!]
  \centering
  \caption{Limiti dei provider di geolocalizzazione IP utilizzati}
  \begin{tabularx}{\linewidth}{L{3cm} L{3.5cm} L{3.5cm} L{3cm}}
    \toprule
    \textbf{Provider} & \textbf{Limite} & \textbf{API Key} & \textbf{Uso} \\
    \midrule
    ip-api.com & 1\,000 req/min & Non richiesta & Primario \\
    ipapi.co   & $\sim$1\,000 req/giorno & Non richiesta & Fallback \\
    \bottomrule
  \end{tabularx}
\end{table}

Per un'applicazione in produzione con traffico elevato, si raccomanda di
valutare provider con piani commerciali (es. MaxMind GeoIP2, ipinfo.io) da
integrare come unica sostituzione in \texttt{app/api/geo/route.ts}.

% ═══════════════════════════════════════════════════════════════════════════════
\section*{Riferimenti}
% ═══════════════════════════════════════════════════════════════════════════════
\addcontentsline{toc}{section}{Riferimenti}
\begin{thebibliography}{9}

\bibitem{w3c-geolocation}
  W3C Geolocation API Specification, \textit{Geolocation API},
  \url{https://www.w3.org/TR/geolocation/}, 2022.

\bibitem{rfc1918}
  Y. Rekhter, B. Moskowitz, D. Karrenberg, G. J. de Groot, E. Lear,
  \textit{Address Allocation for Private Internets}, RFC 1918,
  Internet Engineering Task Force, 1996.

\bibitem{gdpr}
  Parlamento Europeo e Consiglio dell'Unione Europea,
  \textit{Regolamento (UE) 2016/679 — GDPR},
  Gazzetta Ufficiale dell'Unione Europea, 2016.

\bibitem{react-hooks}
  Meta Open Source, \textit{React Hooks Reference},
  \url{https://react.dev/reference/react}, 2024.

\bibitem{nextjs-route-handlers}
  Vercel Inc., \textit{Next.js Route Handlers},
  \url{https://nextjs.org/docs/app/building-your-application/routing/route-handlers}, 2024.

\end{thebibliography}

\end{document}
